\begin{statement}{E.42}{Proved}{A simultaneous extension to every full pro-prime boundary factor away from two and three}
Let $\widehat H$ be the index-$54$ subgroup
of $\widehat B$ determined by the actual
Fermat covering, and for each prime $p$ put
\[
 \mathscr P_p=
 \widehat B\big/
 \ker\bigl(\widehat H\longrightarrow
                    \widehat H^{(p)}\bigr).
 \tag{E.42.1}
\]
The kernel in \textup{(E.42.1)} is characteristic
in $\widehat H$, hence normal in $\widehat B$.
Thus this is the actual quotient extension
\[
 1\longrightarrow\widehat H^{(p)}
   \longrightarrow\mathscr P_p
   \longrightarrow B/H\longrightarrow1.
 \tag{E.42.2}
\]

Fix any of the $432$ multiplier-one source
points of E.24, any full two-adic class-two
fibre of E.41 and any full three-adic
class-two fibre of E.37.
There is one profinite commutator word
realizing all of these prescribed images
and inducing the identity shadow on
every $\mathscr P_p$, for every $p\ge5$.
It is generating in the actual common
source of $\mathscr D$, all the prescribed
two- and three-primary class-two quotients,
and every finite collection of finite
quotients of the groups
\textup{(E.42.1)}.
In particular all prime-power class-two
boundary towers can be retained
simultaneously.
\end{statement}
\begin{proof}
E.37 and E.41 give an actual word
\[
             f\in\widehat H\cap\widehat F'
 \tag{E.42.3}
\]
with the stated two- and three-adic images
and the fixed $\mathscr D$ image.
Here \textup{(E.42.3)} is an exact subgroup
statement: the retained source word lies
in $F'$, every correction lies in
$K_6\cap\widehat F'$, and
$\widehat F'\subset\widehat H$ by the
defining abelian quotient of the Fermat
cover.

For a prime $p$, let $e_p\in\widehat{\mathbb Z}$
have component $1$ at $p$ and $0$ at all
other primes.  The continuous power map
on the closed procyclic subgroup generated
by $f$ gives the explicit replacement
\[
                       f^\ast=f^{\,e_2+e_3}.
 \tag{E.42.4}
\]
It still belongs to the closed subgroup
$\widehat F'$.
Every finite retained source group
$\mathscr D$ has order divisible only by
two and three.  If its image of $f$ has
order $2^a3^b$, then
$e_2+e_3\equiv1\pmod{2^a3^b}$.
Thus \textup{(E.42.4)} preserves that
image exactly.
The continuous power map on any pro-two
or pro-three group only uses the corresponding
component in $\mathbb Z_2$ or $\mathbb Z_3$.
Both components here are $1$, so all
prescribed two- and three-adic images
are unchanged.
On $\widehat H^{(p)}$ for $p\ge5$ the
component of the exponent is $0$.
Consequently
\[
              f^\ast\big|_{\widehat H^{(p)}}=1
                            \qquad(p\ge5).
 \tag{E.42.5}
\]
This uses the full pro-$p$ group, without
passing first to an abelian or class-two
quotient.

The based multiplier-one source formula
with conjugating word $1$ sends the
literal generators $s,r$ to themselves.
By \textup{(E.42.5)}, its evaluation in
each $\mathscr P_p$ is therefore the
identity map of that actual quotient.
Both order relations and generation there
follow from this explicit equality.
The two- and three-primary relations are
those already proved for $f$.
Thus \textup{(E.42.4)} gives compatible
homomorphisms to every indicated finite
common source.

We prove joint generation on those sources,
retaining the actual $H$ maps.
Let $G_{23}$ be any finite common source
of $\mathscr D$ and chosen two- and
three-primary boundary quotients.
The exact inclusion $K_6\subset H$ supplies its
canonical onto map to $B/H$: indeed $N\subset F$,
$N'\subset F'\subset H$, and $N^6\subset F^3\subset H$.
Thus $G_6$, and hence $G_{23}$, has that quotient.
Write
$A_{23}$ for the image of $H$, its kernel.
All prime divisors of $|A_{23}|$ are two
and three.
The shadow on $G_{23}$ is onto by E.41,
and its induced map on $B/H$ is the
identity because its conjugating word
lies in $\widehat H$.
Hence its restriction to $\widehat H$
is onto $A_{23}$: a preimage of an
element in that kernel must itself
have trivial $B/H$ image.

For a finite set of primes $p\ge5$,
choose finite characteristic $p$-group
quotients $A_p$ of $\widehat H^{(p)}$.
Their invariant kernels give actual
quotients $G_p$ of $\mathscr P_p$ with
kernel $A_p$ and quotient $B/H$.
These characteristic quotients are
cofinal, since $\widehat H$ is finitely
generated: intersect the finitely many
automorphic images of any open normal
subgroup.
Repeated quotients at the same prime
may first be combined by taking their
actual common quotient, still a finite
$p$-group in its geometric kernel.

The image of $\widehat H$ in
$A_{23}\times\prod_{p\ge5}A_p$ is the
whole product.  It surjects to every
factor.  A proper subdirect product of
two factors would give nontrivial
isomorphic quotients of those factors,
by taking the two projection kernels.
Their orders would divide two coprime
integers, which is impossible.
Induction gives the displayed full
product, since all the sets of prime
divisors are disjoint.
It follows that the actual joint source
is the fibre product
\[
       G_{23}\mathop{\times}_{B/H}
                         \prod_{p\ge5}G_p,
 \tag{E.42.6}
\]
with geometric kernel
$A_{23}\times\prod A_p$.
Indeed the joint source already surjects
to $B/H$, and the proved equality of
its geometric kernel gives every element
of this fibre product.

Apply the same argument to the image of
the corrected shadow.
Its $\widehat H$ image is onto $A_{23}$
by the preceding paragraph, and is onto
each $A_p$ because the entire map on
$G_p$ is the identity.
The coprime quotient argument gives
the full geometric product.
Its map to $B/H$ is the identity and
onto, so the shadow is onto
\textup{(E.42.6)}.
This proves joint generation without
assuming that the shadow descends to
an automorphism of the fixed Ferrari
source quotient.

All these maps are supplied by the same
word \textup{(E.42.4)}.
Taking the inverse system of finite
characteristic quotients therefore
proves the assertion for every full
pro-$p$ factor with $p\ge5$.
The class-two quotients are among their
actual quotient maps, which proves the
stated simultaneous all-prime comparison.
\end{proof}
